% Generated draft: requires mathematical and source review.
\documentclass[11pt]{article}
\usepackage[margin=25mm]{geometry}
\usepackage{fontspec}
\setmainfont{Latin Modern Roman}
\usepackage{amsmath,amssymb,mathtools,unicode-math}
\setmathfont{Latin Modern Math}
\usepackage{enumitem,needspace,xurl,hyperref}
\setlength{\parindent}{0pt}
\setlength{\parskip}{5pt}
\setlength{\emergencystretch}{3em}
\newcommand{\N}{\mathbb{N}}
\newcommand{\Z}{\mathbb{Z}}
\newcommand{\Q}{\mathbb{Q}}
\newcommand{\R}{\mathbb{R}}
\newcommand{\C}{\mathbb{C}}
\newcommand{\F}{\mathbb{F}}
\newcommand{\A}{\mathbb{A}}
\newcommand{\PP}{\mathbb{P}}
\newcommand{\E}{\mathbb{E}}
\newcommand{\eps}{\varepsilon}
\newcommand{\dd}{\,\mathrm{d}}
\begin{document}
\section*{UnsolvedMath ID: 30001808}
\textbf{Source problem code:} catalogue entry
\subsection{Definitions and mathematical statement}
Let \(M_g\) denote the moduli space of smooth, connected, projective complex curves of genus \(g\), and let \(A_g\) denote the moduli space of principally polarized complex abelian varieties of dimension \(g\). The Torelli map is
\[
\tau_g:M_g\longrightarrow A_g,\qquad [C]\longmapsto [\operatorname{Jac}(C)],
\]
where the Jacobian carries its canonical principal polarization. Write \(h_{\mathrm S}\) for the Siegel metric on \(A_g\), and write \(h_g=\tau_g^*h_{\mathrm S}\) for its pullback tensor.

Let \(M_g^*\) denote the space called the Satake compactification of \(M_g\) in the supplied question. The usual candidate is the closure of the Torelli image in the Satake compactification of \(A_g\); the supplied record does not specify whether a normalization or another precise realization is intended.

For a fixed coefficient field \(\mathbb K\in\{\R,\C\}\) and fixed analytic conventions, let \(H_{(2)}^k(M_g,h_g;\mathbb K)\) denote the degree-\(k\) cohomology of the chosen complex of square-integrable differential forms with square-integrable exterior derivative. This requires specifying a closed domain for the exterior derivative and whether cohomology is reduced. Let \(IH_{\bar m}^k(M_g^*;\mathbb K)\) denote intersection cohomology with middle perversity, whose value at real codimension \(c\) is \(\bar m(c)=\lfloor(c-2)/2\rfloor\).

The extracted Zucker-type question asks whether, after these conventions have been specified, there is a graded isomorphism
\[
H_{(2)}^k(M_g,h_g;\mathbb K)
\;\cong\;
IH_{\bar m}^k(M_g^*;\mathbb K)
\qquad\text{for every integer }k\geq 0.
\]
The question itself gives no explicit genus range; its supplied dated background singles out \(g>3\). Neither the coefficient choice nor the analytic conventions above are prescribed by the source record. Moreover, wherever the differential of the Torelli map is not injective, its pullback tensor is degenerate. Thus the phrase ``Siegel metric pulled back by the Torelli map'' requires an interpretation at such loci before the left-hand side is fully defined. The record also leaves the stack, orbifold, or coarse-space interpretation unspecified. These unresolved qualifications are part of this statement-only draft.
\subsection{Short English statement}
The question asks whether square-integrable cohomology of curve moduli, using the geometry pulled back from the Siegel metric through the Jacobian map, agrees in every degree with middle-perversity intersection cohomology of its Satake compactification. The intended genus range, compactification, and analytic treatment of the pullback tensor require clarification. The supplied curation labels the problem open, particularly highlighting \(g>3\); this draft does not independently verify that status.
\subsection{Sources}
\begin{itemize}
\item Supplied local UnsolvedMath registry record, permanent ID \(30001808\), exact source problem code \texttt{catalogue entry}, title ``Zucker-Type Cohomology for Satake-Compactified Curve Moduli.'' The record attributes the question to the Oberwolfach Reports Open Problems collection, with proposed year 2011; no individual proposer is supplied.
\item \emph{Arithmetic Groups vs. Mapping Class Groups: Similarities, Analogies and Differences} (2011), Oberwolfach Reports, pp.\ 1637--1708, DOI 10.4171/owr/2011/30. Supplied original-report reference: \url{https://doi.org/10.4171/owr/2011/30}.
\item Supplied literature reference labeled ``Hain--Looijenga related conjecture'': \url{https://arxiv.org/abs/0812.3631}. The record presents this as related material and does not supply an exact equivalence with the extracted question.
\item Supplied literature assessment, checked \textbf{2026-08-21}: the report is described as recording no progress on this Siegel-metric analogue for genus greater than three, and the curation reports finding no later general proof resolving the extracted conjecture. This is supplied dated curation, not an independently verified literature assessment.
\item The registry labels the statement \texttt{corrected\_verified} and says its mathematical content was checked against the official report, with notation, context, and question form normalized. These are supplied verification claims; the report and related reference have not been independently consulted for this draft.
\end{itemize}
\end{document}
